\section{Entrance counts and the uniform upper bound}
\label{rw:upper-section}

The physical interval of Section~11 carries one original-height
derivative logarithm $K$. Its regular pieces have at most one crossing.
The three-saddle and weak endpoint estimates alone allow four closing
zeros, so they cannot give the sharp bound. The terminal arguments
below use the sign left by the regular part to bound $K$ by two zeros
on the whole interval. If a small cycle supplies the stronger entrance
sign, the regular part is zero-free and the bound improves to one.
These are the two counts required by Section~\ref{rw:counting}, in the
fixed-field orientation of Section~2.

We first derive the entrance from a small cycle and carry it along
the regular chain. The primary and secondary terminal arguments then
give the whole-interval count and the uniform upper bound.

\subsection{What a small cycle supplies at the common entrance}

Use the fixed-field orientation $\ell\ge0$ of
Section~\ref{rw:geometry}, with the reflection-symmetric thresholds
of Section~\ref{rw:counting}. The coefficient divisions are those
of the original signed auxiliary maps.
The source representative counts the small cycles and determines the
ratio of the two homogeneous generators. To obtain the entrance sign
of $K$, we then evaluate the regular comparison in the original central
height. This change of representation is necessary because a unit
relating closing functions need not preserve their derivative zeros.

\begin{proposition}\label{rw:small-entrance}
After a common neighborhood restriction, suppose $\rho>0$ and
\eqref{rw:counting:forcing} holds. There is at most one counted
cycle with $p<p_\sharp/4$, and it is simple. A counted cycle with
$p<p_\sharp/8$ implies
\begin{equation}\label{rw:upper:small-ratio}
 \tau_1\sigma_1<0,\qquad
 0<-\tau_1/\sigma_1\le p_\sharp^2/300.
\end{equation}
The actual regular comparison at $p_0=p_\sharp/4$ then satisfies
$\operatorname{sign}K(p_0)=-\operatorname{sign}\sigma_1$.
Thus it supplies the precise signed entrance of hypothesis
\textup{(III)} in Section~\ref{rw:counting}.
\end{proposition}
\begin{proof}
The endpoint action bound of Section~\ref{rw:finite:geometry},
with the fixed primary normalization, places every parameter-dominant
cycle beyond $p>\eta/8\ge p_\sharp$; the same bound is recorded
in the physical cover of Section~11. Hence the small cycles here
belong to the two initial energy classes, and are zeros of the one
function constructed in Proposition~\ref{rw:initial-joint-bound}.
In the source class its outgoing action satisfies
\[
 p^2=2e_+\epsilon h,\qquad |e_+-1|<1/1024.
\]
For a cycle with $p<p_\sharp/4$ this gives
$\epsilon h<p_\sharp^2/31$. On the compact energy class the
same inequality follows from $h\le H_*$ after the scale cutoff
is reduced. Let $J_0$ be the initial-energy interval up to that
upper bound; it contains all these cycles, with their multiplicities.
Its lower energy bound is the fixed $h(z_L)>H_{\rm ECT}+2$.
We do not use any ECT assertion near energy zero.

Write the actual source representative as
\[
 F_{\rm src}=\tau_0 f_1+\sigma_0 f_2+f_3,
 \qquad Q=f_2/f_1,\quad g=f_3/f_1.
\]
The compact overlap and the source derivative estimates in Section~6
give on this single initial-energy interval
\[
 Q\ge q_*>0,\qquad Q'\ge q_0>0,\qquad Q\le C_Qh.
\]
On the compact interval these inequalities follow from the first two
ordered Wronskians and fixed coefficient convergence; on the source
tail they follow from $Q=hC/A$, $Q_h>1/5$. The fixed overlap joins
the same two columns, so their signs can be tested on the union.

On $J_0$ we need the line-breaking contribution and its ordinary
energy derivative to be exponentially smaller than the homogeneous
coefficient norm. The signed compact contribution is $\ell$ times a
bounded coefficient. On the source part it is the exact logarithmic
forcing divided by $h^{3/2}A$.
Since $\widehat\ell=\ell/s$ and $s^{-1}=e^{h+1-s}\le e^{h+1}$,
\eqref{rw:counting:forcing} and the source logarithmic estimates
yield
\begin{equation}\label{rw:upper:small-forcing}
 |g|+|g'|\le C\exp\left(1-\frac{11p_\sharp^2}{620\epsilon}\right)
                      (|\tau_0|+|\sigma_0|)
       =:\vartheta_0(|\tau_0|+|\sigma_0|).
\end{equation}
Here $1/20-1/31=11/620$. The same bound applies to the fixed cutoff
terms on the compact overlap. The source logarithms are evaluated at
$\widehat\ell=o(1)$ on this smaller interval, so their denominators
have fixed margins. The derivative is the actual fixed-coefficient
energy derivative, not an independent source derivative.

If both homogeneous generators vanish, forcing gives $\ell=0$,
and the center identity contributes no isolated zero. Otherwise
\[
 F_{\rm src}/f_1=\tau_0+\sigma_0Q+g,
 \qquad (F_{\rm src}/f_1)'=\sigma_0Q'+g'.
\]
If $|\sigma_0|q_0>\vartheta_0(|\tau_0|+|\sigma_0|)$,
this derivative has a strict sign on all of $J_0$. Every zero is
therefore simple and there is at most one. In the complementary
case $|\sigma_0|\le2\vartheta_0|\tau_0|/q_0$ at small scale,
and
\[
 |F_{\rm src}/f_1|\ge |\tau_0|
 \left(1-\frac{2\vartheta_0 C_Qp_\sharp^2}{31q_0\epsilon}
                         -2\vartheta_0\right)>0.
\]
The exponential absorbs $\epsilon^{-1}$, so there is no zero.
This proves uniqueness and simplicity on the whole small-cycle set.

At the smaller cycle, with energy $h_y$, the same equation gives
\[
 -\frac{\tau_0}{\sigma_0}
 =Q(h_y)+g(h_y)/\sigma_0,
 \qquad |g/\sigma_0|\le\vartheta_0(1+|\tau_0/\sigma_0|).
\]
For small $\vartheta_0$, absorb the ratio on the right. The positive
lower bound $Q\ge q_*$ proves the strict negative sign of
$\tau_0/\sigma_0$. Moreover $\epsilon h_y<p_\sharp^2/127$.
The independent source constants give $C/A\to2/5$. Choose a fixed
$H_1\ge H_*+1$ so that its scale-tail error is below $1/100$,
then choose $p_\sharp$ small so that the coefficient error on this
energy range is below $1/100$. For $h\ge H_1$, $Q\le0.42h$;
below it $Q\le C_QH_1$. Therefore
\[
 \epsilon Q(h_y)\le
 \max(0.42p_\sharp^2/127,\epsilon C_QH_1)
                    \le p_\sharp^2/302
\]
after reducing the primary scale. The strict inequality
$0.42/127<1/302$ and the exponential error leave the margin to
$1/300$. Since $\tau_1/\sigma_1=\epsilon\tau_0/\sigma_0$,
this proves \eqref{rw:upper:small-ratio}.

The generator ratio now determines the entrance sign of the original
central comparison.
The physical-fibre construction in Section~11 includes the regular
point $p_0=p_\sharp/4$. The two regular columns in
Lemma~\ref{rw:regular-columns} satisfy
$C/A=p^2/5+O(p^3)$, uniformly on the normalized compact coefficient
set. Thus at this fixed point $C/A\ge p_\sharp^2/85$ after
choosing $p_\sharp$ small. Let
$w=(\tau_1,\sigma_1)/(|\tau_1|+|\sigma_1|)$.
Equation \eqref{rw:upper:small-ratio} makes
$w_1+w_2C/A$ have sign $\sigma_1$ and modulus at least
$p_\sharp^2/200$. The exact regular representation, divided by
this generator norm, is
\[
 K(p_0)/(|\tau_1|+|\sigma_1|)
          =-A(p_0)(w_1+w_2C(p_0)/A(p_0))+e(p_0).
\]
The positive $A(p_0)$ has a fixed margin, and regular coefficient
convergence together with the fixed-field forcing makes $e(p_0)$
smaller than half the preceding margin. Consequently $K(p_0)$ has
sign $-\sigma_1$, as required. All choices preceded the fixed field
and its cycle. Reflection transfers this statement with its actual
outgoing label and the identity \eqref{rw:counting:reflection}.
\end{proof}

\subsection{The complete regular chain}

The regular primary, wedge and secondary comparisons form one chain.
A common crossing direction on their actual overlaps gives a single
zero bound for $K$. Near the parabolic corner, the primary coefficient
cone shows that the outward direction is $-\operatorname{sign}C$,
with $C=\bar d/\lambda$.

\begin{proposition}\label{rw:regular-chain}
Assume the fixed-field inequality \eqref{rw:counting:forcing}.
The regular initial part of the ordered physical interval has at
most one zero of $K$, and that zero is simple. Near $\mathcal P$
this assertion includes the finite connection, the wedge and the
regular secondary tags, with their actual overlaps.

If a primary regular zero has occurred, its subsequent sign and
outward derivative have sign $-\sigma_1$ on the remaining
regular chain. The same conclusion holds if its initial point
has
\begin{equation}\label{rw:merge:strong-entry}
 \tau_1\sigma_1<0,\qquad
 \operatorname{sign}K=-\operatorname{sign}\sigma_1;
\end{equation}
in that case the regular chain is zero-free.
If the first zero instead occurs in the wedge or in the regular
secondary part, its subsequent sign and outward derivative have
sign $-\bar d$, and $\bar a\bar d<0$.
In particular a secondary terminal after that zero receives
$\bar aC<0$ and $K$ of sign $-C$.
\end{proposition}
\begin{proof}
Saturation and coefficient transport give the positive-central
interval with fixed-fibre forcing. The common-input identity
\eqref{rw:counting:input-change} identifies $K$ on its overlaps.

On the primary compact regular part
Lemma~\ref{rw:finite:regular-direction} gives at most one
simple zero and its outward direction $-\sigma_1$.
At a zero, or under \eqref{rw:merge:strong-entry}, the
positive column quotient on a fixed regular overlap gives
\begin{equation}\label{rw:merge:primary-cone}
 |\tau_1|\le C|\sigma_1|,\qquad
 |\nu|\le C e^{-p_\sharp^2/(20\epsilon)}|\sigma_1|.
\end{equation}
The second inequality involves coefficients only, so it persists
as the input moves along the chain.
Near $\mathcal P$, $\bar B+\bar m$ has a fixed positive
margin. Thus $\bar c$ has the sign of $\sigma_1$,
$|\bar a|\le C|\bar c|$, and
$\bar a\bar c<0$: the latter follows from
$(\bar a+\bar c)\bar c<0$.
Equation \eqref{rw:merge:primary-cone} makes
$\nu=o(|\bar c|)$, so $\bar d=\bar c+\nu$ has that same
sign and $\bar a\bar d<0$.
For $\lambda>0$, $C=\bar d/\lambda$ has that sign too.
No division of the error estimate by $\lambda$ was used.

On the wedge the two columns in
\eqref{rw:secondary:wedge-columns}, after the fixed-fibre
forcing column is absorbed, have negative margins and a
strictly monotone quotient in the actual outward direction.
Hence any zero is simple, requires $\bar a\bar d<0$ and
crosses in direction $-\bar d$. The regular secondary
formula \eqref{rw:secondary:regular-limit} has the same
property: its lower constant has sign $-\bar a$, and
$\log(Q/P)$ has sign $C$, with the derivative margin of
Lemma~\ref{rw:secondary:chain-direction}.
Near the center this follows from its divided columns;
away from the center the nonzero value or derivative margin
proves it directly. A zero with $C=0$ would force the
center identity, so does not give an isolated zero.
The strict-focus reciprocal endpoint uses the stated value
or derivative estimate in that same regular secondary
preparation.

Apply Lemma~\ref{rw:counting:merge} on the actual open overlaps.
After a primary zero, the coefficient cone makes $\bar d$ and
$\sigma_1$ have the same sign, so all later crossings have the
same direction. Without a primary zero, all wedge and secondary
crossings have direction $-\bar d$. This gives one zero on the
union; the local direction lemmas keep $K$ departing once it has
sign $-\bar d$.

The finite lower-double connections in
Lemma~\ref{rw:fibre:parabolic-connections} require a separate
observation. On the lower-cone boundary their regular prefix,
finite bridge, wedge and every retained compact regular
secondary interval have one strict action sign, opposite to
$\bar a$. Their lower action dominates on every retained compact
secondary part, so they contribute no zero; a singular end remains
terminal. On the primary prefix
$\tau_1$ has the sign of $\bar a$. If
$\tau_1\sigma_1<0$, the sign $-\sigma_1$ is therefore
opposite to this actual value of $K$. Thus the strengthened
entrance cannot occur in these boundary neighborhoods.

Under \eqref{rw:merge:strong-entry}, the primary part departs from
zero. Its cone \eqref{rw:merge:primary-cone} carries the same direction
to the later pieces, excluding a first zero there. This proves the
strengthened assertion.
\end{proof}

\subsection{Primary terminal intervals with their regular entrances}
\label{rw:anchors:primary}

Let $R$ be the regular prefix of the physical interval and $T$ its
terminal union. On their actual open overlap, regular changes of input
identify the same original-central $K$.

By Proposition~\ref{rw:regular-chain}, $R$ has at most one simple
zero. After a zero, or under the signed entrance below, it supplies
the departure sign $-\sigma_1$ and the coefficient cone
\eqref{rw:merge:primary-cone}. More generally,
Lemma~\ref{rw:finite:regular-direction} gives that derivative sign
whenever $K$ has sign $-\sigma_1$, without requiring the cone;
if $\sigma_1=0$, $K,K_p$ have sign $-\tau_1$ off the center.
A regular zero requires $\tau_1\sigma_1<0$.
These statements hold on the actual overlap, including nonclosing inputs.

\label{rw:anchors:seam-convention}
If $K$ is constant on a terminal interval, a nonzero constant
contributes no zero; a zero constant extends into its regular overlap
and, by the regular count under forcing, belongs to the exact center
case, contributing no isolated closing zero. Otherwise choose a seam
in the overlap away from zeros of $K$ and the terminal derivative
numerator, after a regular zero if one has occurred. The part of $T$ before the seam is then assigned to $R$,
with every multiplicity retained. We fix this seam before classifying
the terminal critical zeros, and use the same convention for the
secondary anchors and the final union.

\begin{lemma}[Entrance counts]\label{rw:anchors:calculus}
On a tail after this seam, a positive-multiplier numerator with at
most $j$ zeros gives $Z(K)\le j+1$, or $Z(K)\le j$ when
$K K'_{\rm out}>0$ at the seam. On an explicitly constructed
comparison continuation, a two-zero bound and opposite nonzero
endpoint signs give at most one zero on every physical subinterval.
If $u$ decreases outward and a positive-multiplier numerator of
$K_u$ strictly increases in $u$, an entrance with $K<0$ and
$K'_{\rm out}<0$ likewise leaves at most one zero, which is simple.
All counts include multiplicity and apply to finite selections at
excluded endpoints.
\end{lemma}
\begin{proof}
Ordinary and away Rolle in Lemma~\ref{rw:counting:rolle} give the
first assertion. Opposite endpoint signs force odd total multiplicity,
so a two-zero continuation has one zero. In the last assertion the
seam numerator is positive and, moving outward, can become negative
at most once. Thus $K$ first decreases below zero and then strictly
increases, allowing at most one simple crossing.
\end{proof}

\begin{proposition}[Regular lower arcs at an upper saddle]
\label{rw:anchors:regular-saddle}
In the chart of Proposition~\ref{rw:finite:regular-saddle}, the terminal
part has at most two zeros of $K$, with multiplicity. If $R$ has a zero,
its remaining terminal part has at most one. The same improvement holds
when $R$ has no zero but the supplied entrance satisfies
\begin{equation}\label{rw:anchors:signed-entry}
 \tau_1\sigma_1<0,\qquad
 \operatorname{sign}K=-\operatorname{sign}\sigma_1.
\end{equation}
Thus $Z(K;R\cup T)\le2$ without a prescribed entrance and it is
at most one under \eqref{rw:anchors:signed-entry} if the preceding
part has no zero.
\end{proposition}
\begin{proof}
Proposition~\ref{rw:finite:regular-saddle} gives the unanchored bound.
In either improved case, the regular overlap supplies
\eqref{rw:anchors:signed-entry} and the cone
\eqref{rw:merge:primary-cone}; coefficient transport supplies
\eqref{rw:finite:saddle-relative-forcing} on this fixed fibre.
Lemma~\ref{rw:finite:saddle-entry-signs} then gives opposite signs
at the seam and $t=0$ for the same original-central $K$ on its
two-zero comparison continuation. Lemma~\ref{rw:anchors:calculus}
leaves one zero there and hence on its physical part. Add the
possible single regular zero.
\end{proof}

At the three-saddle end the derivative numerator may have two zeros.
Ordinary Rolle would then allow three zeros of $K$. We need the signs
forced by the two-critical-zero case, together with the regular entrance,
to remove this extra possibility.

\begin{proposition}[The three-saddle union at a fixed seam]
\label{rw:anchors:three-saddle}
Suppose $T$ is the three-saddle chart of
Proposition~\ref{rw:finite:three-count}, and impose its relative
forcing on the fixed fibre. Then
\[
 Z(K;R\cup T)\le2.
\]
After a zero in $R$, every seam satisfying the preceding conventions
leaves at most one terminal zero. If there is no regular zero and
\eqref{rw:anchors:signed-entry} holds, the union has at most one zero.
All counts include multiplicity and remain valid when central positivity
or a moving lower boundary truncates the passing interval.
\end{proposition}
\begin{proof}
For the physical interval $J=[p_0,p_\gamma]$ of
Proposition~\ref{rw:saturation}, the outer closing orbit supplies the
local forcing in \eqref{rw:finite:three-forcing}.
Indeed, write the actual seam as $p_s=P_s-\psi_+(t_s)$,
with $\psi_+'>0$. The separatrix input $P_s$ has no finite upper
return. Saturation therefore excludes $p_\gamma\ge P_s$, and hence
\[
 p_s\le p_\gamma<P_s,\qquad 0<t_\gamma\le t_s.
\]
The fixed signed connecting maps constructed in
Section~\ref{rw:finite:three-saddle} contain this smaller upper
passage. Its actual lower returns remain on the passing sides
of the two strong separatrices, by
Proposition~\ref{rw:lower-preparation}; thus
$G_+(t_\gamma)>0$ and $G_-(v(t_\gamma))>0$.
Consequently $\gamma$ is an actual closing zero in the same
three-saddle comparison chart. Equation~\eqref{rw:finite:three-forcing}
holds on its fixed coefficient fibre, including the comparison
continuation used below.

Write $u=G_+(t)$, $z=G_-(v(t))$, $v=t^re^w$ as in
Section~\ref{rw:finite:three-saddle}. On the deficit comparison domain,
\[
 M=tuz K_t,\qquad t>0,\quad u>0,\quad z>0.
\]
The outward coordinate is $p=P_s-\psi_+(t)$ with
$\psi_+'>0$. Let $T_s$ be the actual positive-central tail after
the fixed seam $t_s$, and set
\begin{equation}\label{rw:anchors:fixed-seam-number}
 j_s=Z(M;T_s).
\end{equation}
The local bound gives $j_s\le2$ at this fixed seam.
If $j_s\le1$, Lemma~\ref{rw:anchors:calculus} proves the assertions:
regular departure gives $KK_p>0$ after a regular zero or under the
prescribed signed entry.

Suppose now that $j_s=2$. Apply, if needed, the reflection identity to
the whole union, its coefficients, and its already fixed seam so that
$r>1$. It sends $K$ to its negative composed with the corresponding
input change, and changes the signs of $\tau_1,\sigma_1$ together.
It preserves multiplicity, outward ordering and the signed-entry
condition, and transports the already chosen seam and threshold to
their corresponding actual points. The reflected outgoing label is
the original incoming label; the numerical outgoing cutoff need not
be unchanged.

With this local forcing fixed, apply
Proposition~\ref{rw:finite:three-count} on the same deficit comparison
interval, with the original central function $K=\mathcal K$ and
the already fixed seam. Its two-zero classification gives
\begin{equation}\label{rw:anchors:three-end-signs}
 u_0,z_0>0,\qquad \tau_1>0,\qquad
 M(0+)=\epsilon\kappa u_0z_0>0,\qquad K(0+)=-\infty.
\end{equation}
The actual tail's two zeros exhaust the comparison bound.
The same proposition gives $M(t_s)>0$ and $K_p(t_s)<0$ at its
fixed seam.

The regular overlap determines the remaining sign, $K(t_s)$.
For $\sigma_1\ge0$, $\tau_1>0$ and the positive regular columns
give $K(t_s)<0$, with relative forcing absorbed. For $\sigma_1<0$,
a positive value of $K(t_s)$ would have sign $-\sigma_1$, and
Lemma~\ref{rw:finite:regular-direction} would imply $K_p(t_s)>0$,
a contradiction. The seam avoids a zero, so again $K(t_s)<0$.
Together with \eqref{rw:anchors:three-end-signs}, this gives negative
signs at both ends of the continued tail. Ordinary Rolle bounds its
total multiplicity by three, and parity reduces this to two.
Restricting to the physical subinterval preserves that bound.

There cannot have been a regular zero: since $\tau_1>0$, such a zero
would require $\sigma_1<0$ and would leave $K>0$ at the seam,
contrary to the preceding conclusion. Thus the union has at most two
zeros. Finally the no-zero signed-entry alternative also excludes
$j_s=2$. If $\sigma_1>0$, it contradicts $\tau_1\sigma_1<0$;
if $\sigma_1<0$, it would require a positive seam value of $K$.
Only $j_s\le1$ remains, and the away estimate gives one.
\end{proof}

\begin{proposition}[Primary weak thresholds with their regular part]
\label{rw:anchors:primary-weak}
Let $T$ be either a positive-$\bar B$ weak threshold with a lower
root at $P$, as in Proposition~\ref{rw:infinity:weak-loss}, or the weak-center chart of
Proposition~\ref{rw:infinity:weak-center}. Under the actual overlap
and fixed-fibre forcing hypotheses above,
\[
 Z(K;R\cup T)\le2.
\]
After a zero in $R$, the remaining terminal part has at most one zero.
With no regular zero and \eqref{rw:anchors:signed-entry}, the whole
union has at most one. These statements concern the common physical
passing interval, with all whole-path and central positivity conditions.
\end{proposition}
\begin{proof}
The upper input $u$ decreases outward, so each positive-multiplier
numerator of $K_u$ has sign opposite to $K_p$. Throughout this proof
we restrict the quotients to the physical interval where the lower
whole-path test from $\chi(u)$, $u>\max(0,-\gamma)$, the finite
lower deficit and central positivity all hold.

\textit{The weak center.}
Write $d=d_s$ and $s=\varsigma=\bar B+\bar m$ in this proof.
Use the complete numerator \eqref{rw:infinity:weak-center-numerator},
with the actual moving eigenvalue $\gamma$ and the retained term
$-\epsilon F_I(\chi)A_U$ in $C$. On this fixed fibre,
Proposition~\ref{rw:coefficient-transport} supplies
$|\nu|\le\eta(|d|+|s|)$, $\eta\to0$.
With $e_s=e_\varsigma$, Lemma~\ref{rw:infinity:weak-entry-data}
supplies the divided coefficients
\begin{equation}\label{rw:anchors:weak-columns}
 M=d\overline C-t(se_s+\nu e_\nu),\qquad
 e_s>c,\quad \overline C'>1/2,
\end{equation}
on the same coefficient branch, with bounded $C_u^2$ coefficients
and small $\overline C$. Under that forcing the coordinate norm is
comparable to $|\tau_1|+|\sigma_1|$.

If $|s|<\sqrt\eta(|d|+|s|)$, the small-$s$ alternative of that
lemma gives $M/d$ strictly increasing, so $Z(M)\le1$.
Lemma~\ref{rw:anchors:calculus} gives all the assertions in this case.

Otherwise $s\ne0$ and
$e=e_s+(\nu/s)e_\nu>c/2$. If $d=0$, $M=-tse$ has no zero and
the assertions follow by the same Rolle inequalities. If $d\ne0$, set
\[
 R_0=\frac{\overline C}{te},\qquad
 M=dte(R_0-s/d).
\]
Lemma~\ref{rw:infinity:weak-entry-data} gives unimodality of $R_0$,
interval strict superlevel sets, and strictly increasing crossings
of nonpositive levels. In particular,
if $ds<0$, there is at most one zero of $M$, and the preceding
ordinary and away Rolle arguments apply.

It remains that $d$ and $s$ have the same sign. Since $\bar c>0$,
this is the sign of $\sigma_1$. Suppose first that the seam has
$K$ of sign $-\sigma_1$. Its regular derivative has that same sign,
so the seam numerator has sign $d$. Thus the seam lies in
\[
             \{u:R_0(u)>s/d\}.
\]
Moving outward decreases $u$ and can leave this interval at most once.
Consequently the remaining terminal part has at most one numerator
zero. Its away entrance gives at most one zero of $K$. This argument
applies after the possible regular zero and under the prescribed
signed entrance.

If instead the seam value of $K$ has sign $\sigma_1$, no regular
zero has occurred. If $M$ there has sign $d$, the same superlevel
argument gives at most one numerator zero and ordinary Rolle gives
at most two zeros of $K$. If $M$ has the opposite sign, then $K_p$
has the sign of $K$ at the seam; away Rolle and $Z(M)\le2$ again
give at most two. The seam avoids numerator zeros. If the original
overlap point was such a zero, its initial segment was assigned to $R$
before the seam was fixed, so the count still covers the entire union.

\textit{Positive $\bar B$.}
At a simple threshold $\bar a=-\bar c-\bar B/\bar c$ with
$0<\bar B<\bar c^2$, and at the double threshold
$\bar B=\bar c^2$, $\bar a=-2\bar c$. Hence at either threshold
$\tau_1<0$ and $\sigma_1>0$. Lemma~\ref{rw:infinity:weak-entry-data}
gives $K>0$ on the retained compact regular overlap for this same
original central function.

There was therefore no preceding regular zero, since its outward
sign would be $-\sigma_1<0$. For the same reason the prescribed
signed entry cannot occur on this union. It remains to bound the
unanchored terminal continuation of this positive overlap.
Use the numerator and quotient from
Proposition~\ref{rw:infinity:weak-loss},
\[
 M=AC-tE,\qquad E>c,\qquad R_0=C/(tE).
\]
Here $A=\delta_L$ for a simple lower loss, and
$A=\mu_L+\delta_L^2$ for a double loss.
Lemma~\ref{rw:infinity:weak-entry-data} gives $A>0$ and $C>0$
at every numerator zero and interval strict superlevel sets of $R_0$.
If $A\le0$, $M$ has no
zero and the assertion follows. If $A>0$, an entrance with $M>0$
lies in the interval $\{R_0>1/A\}$ and encounters at most one
numerator zero outward; ordinary Rolle gives at most two zeros of $K$.
An entrance with $M<0$ has $K_p>0$ and is away from zero, so
$Z(M)\le2$ and away Rolle give the same bound.

Reflection transports the union, its signs and multiplicities.
The improvement comes from this entrance; an arbitrary weak-chart
subinterval still has the unanchored four-zero closing estimate of
Section~8.
\end{proof}

\subsection{Secondary terminal intervals with their entrances}
\label{rw:anchors:secondary}

Let $R$ be the regular, wedge and regular-secondary prefix, $T$ its
finite-saddle or weak terminal interval, and $O$ their actual open
overlap. We use the common $\mathcal K$ of
\eqref{rw:secondary:complete-log}, the fixed-fibre forcing of
Section~11 and the seam convention of Section~\ref{rw:anchors:primary}.

By Proposition~\ref{rw:regular-chain}, the chain has at most one
simple zero, after which
\begin{equation}\label{rw:secondary:anchor-entry}
 \bar aC<0,\qquad
 \operatorname{sign}\mathcal K=-\operatorname{sign}C,\qquad
 \operatorname{sign}\partial_{\rm out}\mathcal K
                                      =-\operatorname{sign}C .
\end{equation}
Here $C=\bar d/\lambda$ with $\lambda>0$ fixed on the fibre.
Lemma~\ref{rw:secondary:chain-monotone} gives the same departure
from a supplied signed entrance; the positive column quotient on a
compact regular-secondary overlap gives $|\bar a|\le K_1|C|$.

\begin{proposition}[The finite secondary saddle with an entrance]
\label{rw:secondary:saddle-anchor}
Let $T$ be the finite-saddle interval of
Proposition~\ref{rw:secondary:saddle}, with its fixed-fibre relative
forcing. Without a prescribed entrance it has
$Z(\mathcal K;T)\le2$.
If the preceding chain has a zero, then its remaining terminal part
has at most one zero of $\mathcal K$. The same improvement holds
without a preceding zero when $O$ satisfies
\begin{equation}\label{rw:secondary:saddle-signed-entry}
 \bar aC<0,\qquad
 \operatorname{sign}\mathcal K|_O=-\operatorname{sign}C .
\end{equation}
Thus $Z(\mathcal K;R\cup T)\le2$ without a specified entrance,
and the terminal union after an entrance satisfying
\eqref{rw:secondary:saddle-signed-entry} has at most one zero
if its preceding chain has no zero.
All counts include multiplicities and use the actual passing and
positive-central restrictions.
\end{proposition}

\begin{proof}
First separate the lower-double overlay. On a compact regular overlap
before the saddle, the upper logarithms and the relative lower term
are bounded, whereas \eqref{rw:secondary:lower-constant} has arbitrarily
large modulus and sign $-\operatorname{sign}\bar a$ near
$\bar a=\pm2\sqrt{\bar B}$. Under
\eqref{rw:secondary:saddle-signed-entry} this is
$\operatorname{sign}C$, opposite to the entrance sign, so the anchored
case is excluded. The value-domination argument of
Proposition~\ref{rw:secondary:saddle} already gives at most one
zero in the overlay, all crossings having the same nonzero actual-ratio
sign. This count survives restriction by every lower whole-path test.

On the remaining compact regular-lower block the signed comparison
maps extend along the saddle quadrant to $t=0$, with
\begin{equation}\label{rw:secondary:anchor-saddle-log}
 \mathcal K=K_c+\epsilon\kappa\log t+\mathcal E(t),
 \qquad \mathcal E(t)\longrightarrow0 .
\end{equation}
Here $\kappa=r-1$ is the actual ratio difference. The two central outputs
need not both be positive on this whole continuation.

To determine the endpoint sign of this continuation, we need the sign
of the actual ratio difference $\kappa$.
The overlap cone and the transported forcing give
\[
 |\bar a|\le K_1|C|,\qquad
 |\nu|\le\eta_\epsilon(|\bar a|+|C|)
                    \le (K_1+1)\eta_\epsilon|C|.
\]
Near $C=0$, identity \eqref{rw:secondary:saddle-ratio} now reads
\[
 \kappa=-\alpha_s C+\mathcal R_s,\qquad
 |\mathcal R_s|\le
 C_0\bigl(|C|+\epsilon\lambda K_1+
           \epsilon\lambda(K_1+1)\eta_\epsilon\bigr)|C|,
 \qquad \alpha_s>c_0>0 .
\]
Choose the $C$ neighborhood and then the scale restrictions so that
the coefficient of $|C|$ in this error is less than $c_0/2$.
Thus $\operatorname{sign}\kappa=-\operatorname{sign}C$ and
$|\kappa|\ge c_0|C|/2$. Away from this neighborhood the same sign
follows with a strict margin from
\[
 r_0-1=-\frac{2C}{\sqrt{C^2+4M}+C};
\]
the exact divided field is a uniformly small perturbation on the
compact positive-$M$ block. The case $C=0$ is excluded by
$\bar aC<0$. The retained coefficient error $\epsilon\lambda|\nu|$
is therefore small relative to $|C|$, which gives the ratio sign
uniformly even as $C$ tends to zero.

For $|\bar a|>K_0|\kappa|$ the same value-domination argument
already gives at most one zero. In the complementary cone the
unanchored estimate $Z(\mathcal K)\le2$ holds on the entire
saddle comparison continuation. Equation
\eqref{rw:secondary:anchor-saddle-log} gives its endpoint sign
$\operatorname{sign}C$, opposite to the sign at the fixed seam.
The opposite endpoint signs and Lemma~\ref{rw:anchors:calculus}
give one zero on the continuation and at most one on its physical
subinterval. The departure
\eqref{rw:secondary:anchor-entry} keeps the overlap zero-free after a
chain zero; adding that possible zero gives the union count.
\end{proof}

\begin{proposition}[The weak secondary end with an entrance]
\label{rw:secondary:weak-anchor}
Consider the weak row of Proposition~\ref{rw:secondary:real},
first with $C>0$. Without an entrance condition its terminal
interval has $Z(\mathcal K;T)\le2$.
If the preceding chain has a zero, its remaining terminal part
has at most one zero. This improvement also holds when there is
no chain zero but the open overlap satisfies
\begin{equation}\label{rw:secondary:weak-signed-entry}
 \mathcal K|_O<0,\qquad \partial_{\rm out}\mathcal K|_O<0 .
\end{equation}
Consequently the whole chain and terminal union has at most two
zeros of $K$; with the supplied entrance and no chain zero,
it has at most one. Reflection supplies the $C<0$ statement.
These conclusions permit either lower double root to shorten the
passing interval and permit positive-central truncation.
\end{proposition}

\begin{proof}
On the physical fibre Section~11 supplies
\[
 |\nu|\le C\delta_\epsilon(|\bar a|+|\bar d|),\qquad
 \bar d=\lambda C,
\]
from the outer closing zero. After a chain zero, or under the entrance,
$\mathcal K<0$ and $\partial_{\rm out}\mathcal K<0$ there.
The actual outgoing tag is
$P=P_s-\psi(t(u))$ with $\psi'>c$ and $t'>0$,
so outward motion decreases $u$ and
$\operatorname{sign}\partial_{\rm out}\mathcal K
=-\operatorname{sign}\mathcal M$.
Lemma~\ref{rw:anchors:calculus}, with the increasing numerator of
Proposition~\ref{rw:secondary:real}, leaves at most one simple tail
zero on the lower-passing, positive-central interval. Adding the
possible chain zero gives the assertions; reflection supplies $C<0$.
\end{proof}

\subsection{Joining the one terminal union}

We now join the regular departure to the terminal estimates. Each
terminal union includes its constant-sign connections, so the result
counts $K$ on the whole physical interval.

\begin{proposition}[The logarithmic-ratio count on the physical interval]
\label{rw:anchored-unions}
Assume \eqref{rw:counting:forcing}, and let
$J=[p_0,p_\gamma]$, with
$p_0\in\{p_\sharp/8,p_\sharp/4\}$, be the physical interval
of Proposition~\ref{rw:saturation}. Then
\begin{equation}\label{rw:merge:two}
 Z(K;J)\le2.
\end{equation}
If \eqref{rw:merge:strong-entry} holds at $p_0$, then
\begin{equation}\label{rw:merge:one}
 Z(K;J)\le1.
\end{equation}
All zeros are counted with their full analytic multiplicities.
These assertions verify hypothesis \textup{(III)} of
Section~\ref{rw:counting}.
\end{proposition}
\begin{proof}
Write $J=R\cup T$ by Proposition~\ref{rw:ordered-cover}, with at
most one terminal union and every internal constant-sign connection
retained. Section~11 gives its domains and fixed-fibre forcing.
Use the seam and center convention of Section~\ref{rw:anchors:primary}.
If there is no terminal part, Proposition~\ref{rw:regular-chain}
already gives both conclusions.
For each direct derivative row below,
Lemma~\ref{rw:anchors:calculus} gives two tail zeros, or one after
a regular zero or the supplied strong entrance: Proposition
~\ref{rw:regular-chain} gives $KK_p>0$ at that seam.
A direct bound $Z(K;T)\le1$ gives the same union conclusions;
a constant-sign piece contributes no zero.

The following table records the terminal estimates used in the proof.
Every derivative numerator has a positive multiplier on its actual
passing interval.
\begin{center}\small
\begin{tabularx}{\linewidth}{@{}p{.42\linewidth}X@{}}
\toprule
 Terminal description & Estimate or entrance result used\\ \midrule
 Two finite lower roots
       &$Z(K')\le1$; \ref{rw:finite:two-root-count}\\
 One incoming loss at a finite saddle
       &$K<0$ if $r\ge1$, or $Z(K')\le1$ if $r<1$;
                         \ref{rw:finite:upper-loss}\\
 Strict or repeated one-infinite lower contact
       &$Z(K')\le1$; \ref{rw:infinity:strict-lower-loss},
                         \ref{rw:infinity:repeated-lower-loss}\\
 Signed horizontal finite root or cubic end
       &$Z(K')\le1$; signed-horizontal calculation in
          Section~\ref{rw:infinity:signed-horizontal}, or
          Proposition~\ref{rw:infinity:cubic}\\
 Strict focus or nonjoint opposing tail
       &$Z(K)\le1$; \ref{rw:infinity:strict-focus},
                         \ref{rw:infinity:opposing}\\
 Joint contact and separated tail
       &$Z(K')\le1$ on the whole union;
                         \ref{rw:joint:one-critical}\\
 Finite saddle with regular lower legs
       &Anchored result \ref{rw:anchors:regular-saddle}\\
 Primary three-saddle
       &Fixed-seam union result \ref{rw:anchors:three-saddle}\\
 Regular outgoing lower arc at a primary weak upper end
       &$Z(K')\le1$; \ref{rw:infinity:weak-regular}\\
 Primary weak contact or weak center
       &Union result \ref{rw:anchors:primary-weak}\\
 Secondary finite saddle or weak end
       &\ref{rw:secondary:saddle-anchor},
                         \ref{rw:secondary:weak-anchor}\\
 Other secondary end
       &$Z(K)\le1$ or a strict derivative;
                         \ref{rw:secondary:regular},
                         \ref{rw:secondary:real}\\
\bottomrule
\end{tabularx}
\end{center}
The constant-sign lower or upper slow pieces are those of
Lemma~\ref{rw:finite:incoming-contraction} and
Propositions~\ref{rw:infinity:negative-regular},
\ref{rw:infinity:slow-domination}, together with their
reflections and secondary lower-double overlays.
The joint estimate applies to its whole contact--separated union.
At a cubic end the multiplier is $F_+(t)F_-(f(t))>0$ on passage.
It remains to treat interior connections and the entrance-dependent
anchors.

\emph{A lower loss before a finite saddle.}
Suppose an incoming double root precedes a finite upper
saddle. Write its lower parameters as
$\bar B=b_0^2,\bar a=2b_0$ and the upper eigenvalues
as $\alpha_0,-d_0$, with $\alpha_0,d_0>0$.
The root inequality $1/b_0<q_c=1/d_0$ gives
\[
 \tau_1=2b_0+\alpha_0-d_0>b_0+\alpha_0>0 .
\]
The regular action difference is positive at small input
and tends to $+\infty$ at the incoming lower root.
The regular one-simple-zero result therefore makes
$K<0$ throughout the intervening limiting regular
interval. Compact convergence gives its actual sign
on the retained overlaps, and incoming contraction
continues it through the finite connection.
There was no regular zero. The remaining saddle interval
has at most two zeros of $K$ by
Proposition~\ref{rw:finite:upper-loss}; for $r\ge1$
it has none. Hence the complete union has at most two.
If the strong entrance held, $\tau_1>0$ would force
$\sigma_1<0$ and a positive initial value of $K$,
contradicting the negative regular sign.
Thus this configuration cannot occur under that entrance.
The coincident incoming loss needs no intervening
connection and its derivative bound is applied directly.
Reflection proves the opposite orientation with its
corresponding interval and signs.

\emph{A lower root before or after a one-infinite threshold.}
A simple root before the threshold ends passage and
cannot be followed by a second counted part.
For a double outgoing root $r$ strictly before a strict
or repeated upper threshold $P$, the inverse incoming
contraction estimate gives $K>0$ through its finite
connection. The outgoing action diverges while the
incoming action there is regular; the small-input sign
is positive as well, so the regular one-zero result
excludes a preceding regular zero.
For example at the repeated field
$\bar a=-2/r,\bar B=r^{-2},\bar c=2/P,
\bar m=-P^{-2}$, $r<P$ gives
$\tau_1<0$ and $\sigma_1>0$. At a strict threshold
write the positive roots of the upper characteristic
polynomial as $\gamma_+>\gamma_->0$.
The inequality $r<1/\gamma_+$ gives
$\tau_1=\gamma_++\gamma_--2/r<0$ and
$\sigma_1>0$ again.
Thus the strong entrance would require a negative
initial value and is excluded.
The upper contact and separated tail after that connection
have the same strictly negative $K_p$ on their open
overlap: this is the double-before-threshold numerator
in Section~8 and, at a repeated tail, the sign in
\eqref{rw:infinity:opposing-slope}.
Their union has at most one zero. Its preceding positive
connection does not increase this bound.

At a repeated upper threshold with a first lower root
$r>P$, the intervening returning tail is the
constant-sign piece $K<0$ of
Proposition~\ref{rw:infinity:negative-regular}.
The finite-root derivative and its infinite limiting
description are alternative overlapping charts of the
same end, as in Proposition~\ref{rw:ordered-cover}.
For a first simple root
$\bar a=-1/r-\bar Br$ and $\bar Br^2\le1$; hence
\[
 \sigma_1=2P^{-1}(\bar B-P^{-2})<0,\qquad
 \tau_1=2/P-1/r-\bar Br>0.
\]
The same inequalities hold at a double root $r>P$.
The small-input value and the last regular value are
both negative, so again the regular prefix is zero-free.
The strong entrance would have positive sign and is
excluded. The final finite-root piece has a strict
derivative, or the signed horizontal/cubic numerator
has at most one zero. It consequently contributes at
most two zeros, with no count from the preceding
constant-sign connection. Root coincidence with $P$
is one contact case already listed above.

The weak boundary with a lower double strictly before
its threshold uses its constant-sign connection and the
weak-regular numerator of
Proposition~\ref{rw:infinity:weak-regular}.
For an outgoing double, $r<P=1/\bar c$ gives
$\tau_1=\bar c-2/r<0$, $\sigma_1=\bar c/r^2>0$
and a positive initial/connection sign.
For an incoming double, $\bar a=2/r$ gives
$\tau_1=\bar c+2/r>0$ and $\sigma_1=\bar c/r^2>0$;
the incoming action gives the negative connection sign,
and the strong coefficient condition is itself impossible.
In either case the primary
regular prefix has no zero, its value is incompatible
with the strong entrance, and the final numerator has
at most one zero. This gives at most two zeros on the
union. The weak contact exactly at the threshold is
the anchored weak case below.

\emph{Primary anchors.}
Apply the three primary entrance propositions listed in the table.
Coefficient transport supplies their fixed-fibre forcing, including
the primary order-one $\nu$ term; the three-saddle proposition derives
its stronger \eqref{rw:finite:three-forcing} from the same outer cycle.
Their union counts already include the regular prefix and exclude
the incompatible entrances in the two-critical three-saddle and
positive-$\bar B$ weak branches.

\emph{Secondary anchors.}
The complete regular chain supplies
\eqref{rw:secondary:anchor-entry}, either after its zero or from
\eqref{rw:merge:strong-entry}. Together with coefficient transport,
these are the saddle-anchor hypotheses. At the weak end, reflection
to $C>0$ gives $K,K_p<0$, the weak-anchor entrance. The two
propositions give the asserted union counts, retaining lower-double
overlays and their incompatible signs. The other secondary rows use
the table; their regular/focus pieces have already been counted in $R$.

Exhaustion in Section~11 completes the proof: a regular zero leaves
at most one terminal zero; without one the union has at most two.
Under \eqref{rw:merge:strong-entry} the chain is zero-free, each
compatible terminal case leaves at most one zero, and the sign
arguments above exclude the remaining connections.

Physical truncation restricts the counted interval. The endpoint
signs used above belong to the stated comparison continuations.
Under reflection the original central logarithm transforms as
$K_{\rm ref}(q)=-K(p)$ and
$K_{\rm ref}\partial_qK_{\rm ref}=KK_p/q'(p)$.
Thus multiplicities and outward departure are preserved at the
corresponding seam and threshold, while both generator signs reverse.
The reflected outgoing label is the original incoming label.
All noncompact counts use the finite-selection convention of Section~3.
\end{proof}


\subsection{Verification of the four hypotheses}

\begin{theorem}[Uniform upper bound]\label{rw:upper-bound}
There are a fixed open spherical neighborhood $\mathscr U$ of
$\Gamma$ and an open coefficient neighborhood $\mathscr V$ of
$X_0$ in the twelve-dimensional space of quadratic vector fields
such that every field in $\mathscr V$ has total Poincar\'e
multiplicity at most three among its isolated affine periodic
orbits contained in $\mathscr U$.
\end{theorem}
\begin{proof}
First work in the five-parameter family \eqref{rw:family}.
We check the four hypotheses with one common choice of constants.
Choose the independent base-energy threshold and collar of Section~2.
Fix the primitive finite derivative orders before choosing the slow
sections and the larger comparison coefficient boxes of Sections~4--5.
The endpoint preparations of Sections~7--10 and the refined compact
cover of Section~11 then supply finitely many positive margins and
section widths. Choose $p_\sharp\le\eta/8$ smaller than the finitely
many bounds required there and in Proposition~\ref{rw:small-entrance}.
Finally reduce the original coefficient neighborhood so that all
scale cutoffs, relative errors, strict margins and absorptions hold.
These choices precede both the field and any selection of its cycles.
They are made on a reflection-symmetric neighborhood.

Hypothesis \textup{(I)} is Proposition~\ref{rw:initial-joint-bound}
together with the endpoint action bound of Section~7: the two initial
classes have one joint bound two, and every parameter-dominant cycle
has $p>\eta/8$. The branch $\rho=0$ belongs entirely to the initial
classes. No scale quotient is used on that branch.

For $\rho>0$ choose the fixed-field orientation $\ell\ge0$.
Proposition~\ref{rw:outer-forcing} gives the original coefficient
inequality \eqref{rw:counting:forcing} whenever its single monotone
alternative does not already settle the remaining interval.
Proposition~\ref{rw:saturation} then supplies the entire physical
segment from either prescribed entrance to any selected outer cycle.
It includes every intermediate nonperiodic comparison point and
extends to neighborhoods of its regular endpoints. The central
section construction of Section~2 identifies closing-zero orders
with Poincar\'e multiplicities. This verifies \textup{(II)}.

Propositions~\ref{rw:ordered-cover} and
\ref{rw:coefficient-transport} put those same segments in the local
comparison domains with the same fixed coefficient inequality.
Proposition~\ref{rw:anchored-unions} bounds $K$ by two zeros on the
whole union, and by one when the specified coefficient signs and
entrance sign hold. This verifies \textup{(III)}. The count on this
segment already includes the initial-class cycles lying there.
Proposition~\ref{rw:small-entrance} proves the uniqueness, simplicity,
coefficient ratio and regular-section sign required by \textup{(IV)}.

Proposition~\ref{rw:conditional-three} now applies to every finite
selection of counted cycles. Its proof adds actual outer and small
witnesses when needed and does not assume that an outermost cycle
among all cycles exists. Thus every finite selection has total
multiplicity at most three. A larger or infinite isolated collection
would contain a finite selection violating that bound.
Identity returns contribute no isolated cycles.

The exact reflection identities in Section~3 transfer the result
back to the original orientation, using $p_{\rm ref}=q$ and the
inverse return germ. They preserve multiplicity and the actual
entrance condition; an unchanged original $p$ threshold is not
required. Finally Proposition~\ref{rw:slice} supplies one smaller
fixed spherical collar and one open neighborhood in the full
coefficient space. Its near-identity affine conjugacies and positive
time changes preserve isolated-cycle multiplicity. The resulting
$\mathscr U,\mathscr V$ prove the theorem.
\end{proof}
